% ===================================================================
%  اللوحة نمرة ١٦ — المحل الكبير. --- البازار. اللعب.
%
%  Traduction, faite en 2026, en arabe egyptien ecrit, sur l'ido de
%  texto/io. Regles de la colonne : voir 10-lawha-01.tex. Une seule
%  numerotation. Dernier tableau du livret.
%
%  LE THEATRE DE MARIONNETTES (41) EST مسرح العرايس, ET LE PERSONNAGE
%  QUI Y JOUE (43) EST الأراجوز. Ce n'est pas une traduction de
%  Polichinelle : l'Aragoz est la marionnette egyptienne, celle des
%  moulids et des rues, et le mot vient du turc kara göz, « œil
%  noir ». Le livret de 1926 met un Polichinelle en scene ; l'egyptien
%  a le sien, avec son nom et son theatre.
%
%  LA MENAGERIE DE CARTON REPREND LA REGLE PHONOLOGIQUE UNE DERNIERE
%  FOIS, ET SUR DEUX BETES QUE TOUT LE MONDE NOMME :
%
%    تعلب (11)   le renard   MSA : ثعلب
%    ديب (21)    le loup     MSA : ذئب
%
%  Et le levrier (14) est السلوقي, le saluki, qui est le chien de
%  chasse arabe : le cantonais avait mis 格力狗, « greyhound » venu
%  des courses de Hong Kong, la ou l'egyptien avait le sien depuis
%  toujours.
%
%  LE MAGASIN, LUI, FINIT COMME IL AVAIT COMMENCE, EN ITALIEN ET EN
%  TURC :
%
%    شيفونيرة (68)  la commode      de « chiffonnier »
%    جزدان (88)     le porte-monnaie
%    فانوس سحري (2) la lanterne magique
%    شاويش (28)     le sergent      du turc çavuş
%    نشّال (102)     le pickpocket
%    بلبل            le rossignol
%    وطواط (35)     la chauve-souris
%    أصلة (36)      le boa
%
%  ECART DECLARE : le bloc t16-15-1 rend la lettre (d) des deux
%  hemispheres, que le fac-simile ido compose « d) » sans parenthese
%  ouvrante alors que la planche la grave.
%
%  CE TABLEAU CLOT LA COLONNE. Seize planches, 683 blocs, et pas une
%  seule inversion de renvoi due a la langue : l'annexion arabe place
%  le possede avant le possesseur, exactement comme l'ido, et les
%  quinze reprises qu'il a fallu ailleurs etaient toutes des
%  maladresses de redaction, jamais des contraintes de syntaxe. Le
%  chinois avait coute trente-cinq inversions au seul debut de la
%  colonne cantonaise.
% ===================================================================

%%K t16-apar-1 apar
\VUsaut{4.0mm}
\VUtitre{15.0pt}{60}{اللوحة نمرة ١٦}
\VUsaut{2.0mm}
\VUpk{13.2pt}{40}{المحل الكبير. --- البازار.}
\VUpk{13.2pt}{40}{اللعب.}
\VUsaut{3.4mm}

%%K t16-01-1 p
1. --- السنة الجديدة قرّبت. والمحلات الكبيرة حضّرت فاترينات \VUgras{اللعب}
\textsuperscript{(1)} وهدايا رأس السنة.

%%K t16-01-2 p
وأنا طلبت من جدو إنه ياخدني للبازار علشان أشوف المعرض الحلو ده، ووافق.

%%K t16-02-1 p
2. --- الأول وقفنا قدام دروع سلاح حلوة، فيها \VUgras{بندقية}،
و\VUgras{كاسكيت}، و\VUgras{سيف}، و\VUgras{حزام سيف}، وجوز \VUgras{رتب}؛ أو
\VUgras{خوذة}، و\VUgras{درع} \textsuperscript{(3)}، و\VUgras{مسدّس} \textsuperscript{(4)}.
وده كله شدّني أكتر بكتير من \VUgras{الفانوس السحري} \textsuperscript{(2)}.

%%K t16-tit-1 sub
\VUsaut{3.0mm}
\VUpk{10.2pt}{40}{مناظر حلوة.}
\VUsaut{2.6mm}

%%K t16-03-1 p
3. --- وفي نفس القسم كان في \VUgras{حارس} \textsuperscript{(5)} في \VUgras{كشكه}؛
وشكله كأنه بيحرس حديقة حيوان كاملة من الكرتون. وكام حيوان كان هناك:
\VUgras{ببغا} \textsuperscript{(6)} على عصايته، و\VUgras{وحيد قرن} \textsuperscript{(7)}
بقرن على مناخيره، و\VUgras{رنّة} \textsuperscript{(8)}، و\VUgras{نعامة}
\textsuperscript{(9)}، و\VUgras{قرد} \textsuperscript{(10)} بيكسّر عين جمل،
و\VUgras{تعلب} \textsuperscript{(11)} شايل فرخة، و\VUgras{تمساح} \textsuperscript{(12)}
فاتح فكّين ضخمين، و\VUgras{جمل} \textsuperscript{(13)}، و\VUgras{سلوقي}
\textsuperscript{(14)} شيك، و\VUgras{نمر} \textsuperscript{(15)}، وبعدين اتنين حيوانات ما
كنتش أعرفهم واللي هما، على ما بيقولوا، \VUgras{ابن عرس} \textsuperscript{(16)}
و\VUgras{ضبع} \textsuperscript{(17)}. وكان هناك كمان \VUgras{فهد} \textsuperscript{(18)}
و\VUgras{ديب} \textsuperscript{(21)}؛ وقريّب أوي كان في \VUgras{فيران}
\textsuperscript{(19)} صغيّرة لطيفة و\VUgras{فيران} \textsuperscript{(20)} صغيّرة أوي من
كاوتش. وأخيرًا \VUgras{فرس النهر} \textsuperscript{(22)} و

%%K t16-03-1 p suite
\VUgras{الجمل أبو سنام} \textsuperscript{(23)} كمّلوا المجموعة. وكنت هفضل هناك ساعات
أكتر لولا إن جنبها كان في معسكر رائع مليان \VUgras{عساكر رصاص} \textsuperscript{(29)}
اللي شدّ انتباهي.

%%K t16-tit-2 sub
\VUsaut{4.0mm}
\VUpk{10.2pt}{40}{عساكر وحرب.}
\VUsaut{2.8mm}

%%K t16-04-1 p
4. --- وجدو، اللي \VUgras{مصاب} \textsuperscript{(24)} واللي اضطرّ يسيب الجيش بسبب
\VUgras{جرح} خلّاه ياخد \VUgras{نيشان الحرب}، بيحبّ أوي يحكي لي \VUgras{الحملات}
اللي اشترك فيها. وكان مبسوط وهو بيشرح حركات الجيش الصغيّر ده المختلفة. ومجموعة
عساكر حقيقيين كانوا بيزوروا البازار: \VUgras{عسكري مهندسين} \textsuperscript{(25)}،
و\VUgras{صيّاد جبال} \textsuperscript{(26)} لابس \VUgras{بيريه}،
و\VUgras{باشجاويش} \textsuperscript{(27)} مشاة، و\VUgras{شاويش مدفعية}
\textsuperscript{(28)} بـ\VUgras{شريط} دهب على كمّه، وقفوا جنبنا يسمعوا الشرح.

%%K t16-05-1 p
5. --- وجدو، وهو فخور أوي إن عنده سامعين مهمّين كده، ارتجل خطة معركة كاملة.

%%K t16-05-2 p
المدينة المحصّنة، بـ\VUgras{أسوارها} و\VUgras{خنادقها}، كان \VUgras{جيش العدو}
محاصرها، واللي بعد \VUgras{ضرب} طويل كان مستعدّ \VUgras{يهجم}. بس
\VUgras{المحاصرين}، وهما مضطرّين من الجوع، جرّبوا \VUgras{يخرجوا} على
\VUgras{معسكر} \VUgras{المحاصِرين}. وبعد ما ردّوا \VUgras{الهجوم} بـ\VUgras{قتال}
عنيد حوالين \VUgras{المعسكر}، \VUgras{المنتصرين} من المحاصِرين أطلقوا
\VUgras{خيّالتهم} علشان يطاردوا المهاجمين \VUgras{المهزومين}. ودول
\VUgras{انسحبوا} وهما بيغطّوا \VUgras{ميدان القتال} بـ\VUgras{الموتى}
و\VUgras{الجرحى}. و\VUgras{حمّالين النقّالة} شالوا دول لـ\VUgras{نقطة الإسعاف}
علشان \VUgras{الجرّاح} يعالجهم.

%%K t16-05-3 p
وبالرغم من المقاومة البطولية بتاعة كام \VUgras{أورطة} اتصفّوا في \VUgras{مربّع}،
\VUgras{الهزيمة} كانت

%%K t16-05-3 p suite
كاملة، وباقي الجيش \VUgras{هرب}، وساب \VUgras{أسرى} كتير. والعدو راح يكمّل
\VUgras{نصره} ويحتلّ المدينة. ويمكن ده يبقى آخر الحملة، وبعد \VUgras{هدنة}، يبقى
ممكن يوقّعوا \VUgras{معاهدة صلح}.

%%K t16-tit-3 sub
\VUsaut{4.4mm}
\VUpk{10.2pt}{40}{هدايا رأس السنة.}
\VUsaut{2.8mm}

%%K t16-06-1 p
6. --- والعساكر الصغيّرين، اللي كانوا بيضحكوا وهما بيسمعوا حكاية المحارب القديم
الحلوة، سألوه كام سؤال تاني. وأنا لفّيت حوالين الفرشة علشان أبصّ على
\VUgras{قوس آلي} \textsuperscript{(32)} حلو و\VUgras{قوس} \textsuperscript{(33)}
بـ\VUgras{سهمه} \textsuperscript{(31)}. وهناك لقيت مجموعة حيوانات تانية: اتنين
\VUgras{طيور مغنّية} \textsuperscript{(34)}: \VUgras{بلبل} آلي و\VUgras{هزار}
بيغنّي بكل حلقه، وسلسلة حيوانات غريبة: \VUgras{وطواط} \textsuperscript{(35)}،
و\VUgras{أصلة} \textsuperscript{(36)}، و\VUgras{سلحفاة} \textsuperscript{(37)}،
و\VUgras{عجل بحر} \textsuperscript{(38)}، و\VUgras{حوت} \textsuperscript{(39)}،
و\VUgras{قرش} \textsuperscript{(40)} وشّه وحش.

%%K t16-07-1 p
7. --- وما وقفتش كتير أتأمّلهم، علشان شفت اللي كنت بحلم بيه: \VUgras{مسرح عرايس}
\textsuperscript{(41)} حلو أوي بـ\VUgras{ديكور} رائع و\VUgras{ستارة} حمرا تسحر.
وعلى \VUgras{المسرح}، اتنين ممثّلين كأنهم بيلعبوا \VUgras{دور} في
\VUgras{مسرحية} مشوّقة أوي، علشان \VUgras{المتفرّجين} الصغيّرين من الكرتون،
المركّبين في اللوچات أو القاعدين على \VUgras{الفوتيهات} وفي \VUgras{الصالة} ورا
\VUgras{قائد الأوركسترا}، كانوا بيصفّقوا بحماس.

%%K t16-07-2 p
وللأسف، المسرح ده كان غالي أوي.

%%K t16-08-1 p
8. --- وكمان اختيار اللعب كان صعب، علشان في الفاترينات كانت متكوّمة كل أنواع
اللعب: \VUgras{أرلكان} \textsuperscript{(42)}، و\VUgras{أراجوز} \textsuperscript{(43)}،
و\VUgras{بيّارو} \textsuperscript{(44)}، و\VUgras{بوم} \textsuperscript{(45)} بيرفرف
بجناحه، و\VUgras{شحرور} \textsuperscript{(46)} بيصفّر، و\VUgras{كلاب بودل} بتنبح،
و\VUgras{فونوغراف} \textsuperscript{(47)}، و\VUgras{لعب الصبر}. وواحدة من اللعب دي
سلّتني أوي: هي

%%K t16-noto-1 noto
\VUnotes{143.2mm}{%
\textit{(*) بُصّ اللوحة نمرة ٦.}%
}

%%K t16-08-1 p suite
كانت بتمثّل \VUgras{معركة بحرية} \textsuperscript{(48)}، اللي فيها \VUgras{مركب} بيهجم
عليها \VUgras{قراصنة} جنب بلد مزروعة \VUgras{شجر جوز هند}.

%%K t16-09-1 p
9. --- وقدرت أتأمّل العجايب دي كلها بسهولة، علشان ما كانش في غير ناس قليلة في
المحل، بالكاد كام واحد بيتفرّج، و\VUgras{دراجون} \textsuperscript{(49)}،
و\VUgras{محصّل} \textsuperscript{(50)} كان بيقدّم \VUgras{كمبيالة} في \VUgras{الخزنة}
\textsuperscript{(51)} الرئيسية.

%%K t16-10-1 p
10. --- وسألت جدو عن فايدة الدفاتر المحطوطة على ترابيزات ورا \VUgras{الكاشيرة}؛
فردّ عليّ إنها \VUgras{دفاتر الحسابات}.

%%K t16-tit-4 sub
\VUsaut{3.6mm}
\VUpk{10.2pt}{40}{الوقوف قدام الفاترينات.}
\VUsaut{2.8mm}

%%K t16-11-1 p
11. --- ولما سبنا قسم اللعب، رحنا قسم السروج نبصّ على \VUgras{السروج}
\textsuperscript{(53)}، و\VUgras{الركايب} \textsuperscript{(54)}، و\VUgras{اللجم}
\textsuperscript{(55)}، و\VUgras{الشكايم} \textsuperscript{(56)}، و\VUgras{الكرابيج}.
ولما \VUgras{رئيس القسم} \textsuperscript{(57)} كان بيدّي كام معلومة لجدو، رحت أبصّ
على فاترينة \VUgras{البورسلين} و\VUgras{الكريستال} \textsuperscript{(58)} اللي فيها
\VUgras{سلطانيات سلطة} حلوة و\VUgras{برادات} \textsuperscript{(59)} رقيقة.

%%K t16-12-1 p
12. --- وعلشان نطلع الدور الأول، عدّينا ورا فاترينة \VUgras{الكورسيهات}
\textsuperscript{(60)}، جنب محل الشرّابات، اللي كانت معروضة فيه
\VUgras{لباسات} \textsuperscript{(63)} حلوة أوي. وقريّب، \VUgras{بايعة}
\textsuperscript{(61)} كانت بتخلّي \VUgras{زبونة} \textsuperscript{(62)} تختار
\VUgras{حرير} \textsuperscript{(64)} حلو.

%%K t16-12-2 p
ولما كنت بطلع السلّم، لاحظت إن المحل مزيّن بـ\VUgras{فوانيس} \textsuperscript{(65)}
يابانية، و\VUgras{شماسي} \textsuperscript{(66)}، و\VUgras{نخل} \textsuperscript{(70)}
ضخم.

%%K t16-13-1 p
13. --- وفي الدور الأول، ما كانش في حاجة كتير تتشاف؛ بس عفش (*)،
و\VUgras{خزن} \textsuperscript{(67)}، و\VUgras{شيفونيرات} \textsuperscript{(68)}،

%%K t16-13-1 p suite
و\VUgras{عربيات أطفال} \textsuperscript{(69)}. بس منظر المحل كله كان \VUgras{مسلّي}
أوي.

%%K t16-14-1 p
14. --- وتحت رجلينا، شفت الآلات الموسيقية: \VUgras{قيثارات} \textsuperscript{(71)}،
و\VUgras{جيتارات} \textsuperscript{(72)}، و\VUgras{مندولين} \textsuperscript{(73)}،
و\VUgras{كورنيت} \textsuperscript{(74)}، و\VUgras{كلارينيت} \textsuperscript{(75)}،
وكمنجات بـ\VUgras{أقواسها} \textsuperscript{(76)}، وحتى \VUgras{طبلة كبيرة}
\textsuperscript{(77)}. وأبعد شوية، في قسم الورق، \VUgras{طالب} \textsuperscript{(78)}
كان بيفاصل \VUgras{البايع} \textsuperscript{(79)} على علبة كراريس. وجنبهم، ميّزت
\VUgras{كتاب ألفباء} \textsuperscript{(80)} حلو.

%%K t16-14-2 p
وفي نصّ المحل كانت متنصّبة \VUgras{شجرة كريسماس} \textsuperscript{(81)} عالية مليانة
شموع، وحاجات صغيّرة، وكل أنواع اللعب: \VUgras{خيل خشب} \textsuperscript{(82)}،
و\VUgras{عرايس} \textsuperscript{(83)}، وهكذا.

%%K t16-14-3 p
و\VUgras{قسم العطور} \textsuperscript{(84)} بـ\VUgras{معجون السنان} وبـ\VUgras{العطور}
المختلفة كان بينشر ريحة حلوة أوي وصلت لنا.

%%K t16-tit-5 sub
\VUsaut{3.4mm}
\VUpk{10.2pt}{40}{بنشتري حاجة.}
\VUsaut{2.8mm}

%%K t16-15-1 p
15. --- وعلشان جدو ابتدا يشوف إن الوقت طال، نزلنا من السلّم الكبير. ووقف شوية
قدام \VUgras{لوحة فلك} \textsuperscript{(85)} بتمثّل، على ما قال لي،
\VUgras{مذنّبات} \textsuperscript{(a)}، و\VUgras{خسوف القمر وكسوف الشمس}
\textsuperscript{(b)}، ووردة رياح بـ\VUgras{الجهات الأربعة الأصلية} \textsuperscript{(c)}
\VUgras{(الشمال، والجنوب، والشرق، والغرب)}، وخريطة \VUgras{نصفي الكرة}
\textsuperscript{(d)}، و\VUgras{الكواكب} \textsuperscript{(e)}، و\VUgras{المجموعة
الشمسية} \textsuperscript{(f)}. وأنا ما فهمتش حاجة من ده كله، وكان اللي شدّني أكتر
بكتير هو بدلة \VUgras{صبي التوصيل} \textsuperscript{(86)} المزركشة اللي عدّى،
و\VUgras{المراوح} \textsuperscript{(87)} الحلوة اللي كانت جنبي، قريّب من
\VUgras{الجزادين} \textsuperscript{(88)} و\VUgras{شنط الأوراق} \textsuperscript{(89)}.

%%K t16-16-1 p
16. --- واتفرّجت كمان على الفرشة على \VUgras{الريش} \textsuperscript{(90)}،
و\VUgras{أبازيم الحزام} \textsuperscript{(91)}، و\VUgras{الورد الصناعي}
\textsuperscript{(94)} الحلو أوي المرتّب بشكل لطيف في مقاطف. و\VUgras{البنفسج}،
و\VUgras{المنثور}،

%%K t16-16-1 p suite
و\VUgras{الصفير} \textsuperscript{(95)}، و\VUgras{الجيرانيوم} \textsuperscript{(96)}،
و\VUgras{الزنبق} \textsuperscript{(97)}، و\VUgras{أبو خنجر} \textsuperscript{(98)} كانوا
متقلّدين كويس أوي.

%%K t16-tit-6 sub
\VUsaut{4.4mm}
\VUpk{10.2pt}{40}{هدية جدو.}
\VUsaut{2.8mm}

%%K t16-17-1 p
17. --- وطلبت من جدو يشتري لي واحدة من \VUgras{فرش السنان} \textsuperscript{(92)}
اللي كانت في علبة جنب \VUgras{بنس الشعر} \textsuperscript{(93)}. ووافق، ورحت بنفسي
أدفع مشترياتي عند الخزنة، اللي جنبها كانوا بيحضّروا \VUgras{طرد}
\textsuperscript{(100)} كبير لـ\VUgras{زبون} \textsuperscript{(99)}.

%%K t16-18-1 p
18. --- وفجأة حصلت لخبطة خفيفة بين الناس الواقفين جنب \VUgras{التماثيل البرونز}
\textsuperscript{(101)} الفنّية. \VUgras{نشّال} \textsuperscript{(102)} اتقبض عليه
متلبّس من \VUgras{مفتّش} \textsuperscript{(103)}، وهو بيحاول يسرق واحد. واهتمّوا
إنهم يجيبوا عسكري بوليس علشان \VUgras{يقبض} عليه، وأول ما خرجنا، كانوا بيودّوا
الجاني للسجن.

\VUsaut{6.0mm}
\VUfilet{22mm}
